% mage-auto-heading: foundation-models-the-source-of-commodity-intelligence
% mage-auto-heading: four-properties-of-commodity-intelligence
% Authoritative LaTeX manuscript; Markdown is a generated compatibility projection.

% <!-- part-title: The New Engineering Problem -->

% <!-- chapter-title: Agentic Machinery and Its Properties -->

\chapter{Agentic Machinery and Its Properties}

% <!-- section-terms: engineering-substrate, foundation-model, agentic-harness -->

% <!-- point: engineering-methods-fit-the-machinery | The roles of the model, harness, and environment explain where engineers can exercise control. | terms: engineering-substrate, foundation-model, agentic-harness -->

Engineering methods must fit the machinery they govern. \autoref{agentic-substrate} distinguishes three components of an agentic engineering system: the foundation model that supplies general capabilities, the harness that turns those capabilities into actions, and the engineered environment in which agents work. Understanding their respective roles helps explain where engineers can exercise control.

\begin{figure}[htbp]
\centering
\includegraphics[width=\linewidth]{book/latex/assets/agentic-substrate-f28785289f.pdf}
\caption{\emph{Commodity intelligence in a governed engineering environment.} General capabilities from a foundation model become useful engineering machinery when connected to an agentic harness and a governed engineering environment. The environment contains both the existing system and the engineering structure used to guide and govern work. Engineers develop and maintain this environment.}
\label{agentic-substrate}
\end{figure}

\section{Foundation Models: The Source of Commodity Intelligence}\label{heading:foundation-models-the-source-of-commodity-intelligence}

% <!-- section-terms: engineering-substrate, foundation-model, agentic-harness, engineered-environment -->

% <!-- point: model-plus-harness-form-a-new-engineering-substrate | Coding agents combine capabilities acquired through foundation-model training with harness facilities that turn responses into action; engineers can shape their operating environment without changing trained weights. | terms: engineering-substrate, foundation-model, agentic-harness, engineered-environment -->

% <!-- index-def: engineering-substrate -->

The coding agents available in 2026 are built primarily around large language models (LLMs), which serve as their \textbf{foundation models}. These models acquire general capabilities through training on large collections of text, software, and other data.\mageaside{Engineering predates machine learning in its use of models as purposeful representations of systems. Machine learning later adopted the same word for a learned approximation of a function---also reasonably a model, but a different kind. This book gives the engineering usage precedence: unqualified model means a purposeful representation of an engineered system or some property of it. I therefore use foundation model for the trained machine-learning system and agent for the acting system built around one through a harness, tools, and working state. Where the distinction matters, I will not ask the reader to infer which ``model'' I mean.} Publicly available source code, including open-source repositories, has been an important source of training material for coding models~\citep{chen2021codex}. Training mixtures can also include licensed material, and their full composition is not always disclosed~\citep{openai2023gpt4}.

The goal of training is to develop a model that can make useful predictions and generalize beyond its training examples. Training typically begins with a prediction objective: given a sequence of tokens, predict the next one. Tokens are units of representation, often words, parts of words, or punctuation. Training adjusts numerical parameters, or weights, to improve these predictions. Additional training can improve how a foundation model follows instructions and responds to human requests~\citep{ouyang2022instructions}. At inference time, the trained model processes the context supplied to it and generates a response. An \emph{agentic harness} turns some of those responses into actions, such as searching a repository, modifying files, running tests, or invoking development tools. Tool results can then be returned as context for further reasoning~\citep{yao2023react}.

A foundation model trained on public software may know programming languages, libraries, common architectures, and established engineering practices. It ordinarily will not know the current state of the particular system an engineer asks it to modify. That system may be proprietary, may have changed since training, or may contain local design decisions that general programming knowledge cannot supply. Even a seemingly local change may depend on an architectural boundary established years earlier, a configuration convention shared across services, or an operational constraint that appears nowhere near the code being edited. Agents must recover relevant knowledge from the codebase, documentation, configuration, and other available evidence. General programming knowledge must be connected to the particular system and the obligations its changes must preserve.

Engineers can influence a language-based agent in two broad ways. They can modify the foundation model through training or fine-tuning, changing its weights. Alternatively, they can change what the agent receives and the environment in which it operates: instructions, system knowledge, tools, permissions, and feedback. The agentic harness provides some of the machinery for exercising this control, including context management, tool access, and hooks through which engineers can observe or constrain actions. Training a frontier foundation model requires resources far beyond those of an ordinary software engineering project. Fine-tuning is more accessible, particularly through techniques that update only a fraction of the model's parameters~\citep{hu2022lora}, but still requires suitable training data, evaluation, and maintenance. It may also be unavailable for a particular foundation model. By contrast, engineers can often configure the harness and establish project-specific controls without modifying the model's weights.

MAGE concentrates on this second approach. The engineer acquires general-purpose intelligence as a service or software component, then engineers the environment in which that intelligence performs consequential work. The foundation model supplies general capabilities, the harness turns its outputs into actions and exposes points of control, and the governed engineering environment supplies the knowledge, obligations, and mechanisms governing those actions.

% <!-- box-family: inset -->

% <!-- inset-domain: REASONING -->

% mage-box-prefix: Inset —

\begin{bookbox}{inset}{Do LLMs Reason?}
\label{do-llms-reason}
Whether LLMs genuinely reason remains a subject of scientific debate. Bender and colleagues challenged the identification of linguistic fluency with understanding~\citep{bender2021stochastic}, while Bubeck and colleagues demonstrated broad problem-solving behavior in an early GPT-4~\citep{bubeck2023sparks}. More controlled experiments have found evidence of learned internal representations: studies of Othello-playing models recovered representations of board state and demonstrated their influence on predictions~\citep{othello-gpt,nanda2023othello}, while Jin and Rinard found evidence that models trained on programs had learned aspects of program semantics~\citep{jin2024semantics}. These studies establish different things: observable problem-solving behavior, learned structure, and evidence connecting internal representations to computation. None establishes that an agent will reason reliably about arbitrary production software. For this book, reasoning is a functional capability: using available information to draw inferences, evaluate alternatives, and work toward a solution. LLM-based systems can supply enough of that capability to undertake substantial engineering work. The question for engineering is how to put it to work under informed control.
\end{bookbox}

\section{Four Properties of Commodity Intelligence}\label{heading:four-properties-of-commodity-intelligence}

% <!-- point: four-foundation-model-properties | Four properties of contemporary foundation models matter to engineering: broad semantic reasoning supports delegated implementation, probabilistic output calls for independent assurance, bounded working state motivates external representations, and inexpensive repeated reasoning increases the need for scalable governance. | terms: foundation-model -->

Four properties of contemporary foundation models matter to the engineering argument.

\textbf{Broad semantic reasoning.} A foundation model can interpret instructions, synthesize information, draw inferences, and search for solutions across open-ended domains. Give it an abstraction---a schema, state table, architectural model, or policy---and it can work with that abstraction rather than merely execute a fixed script. This makes it possible to delegate implementation decisions that engineers have not specified in advance. A plausible decision may nevertheless be inappropriate for the system at hand. An agent may select a familiar architecture or design pattern that works well elsewhere but violates a local constraint, introduces unnecessary complexity, or conflicts with an established design decision.

\textbf{Probabilistic output.} Repeated invocations may produce different implementations or judgments, and its confidence in its own result is not independent evidence of correctness. An agent may also fail consistently when its attempts share an incorrect assumption. Asking another foundation model to review the work can improve reliability, but the errors of the producer and evaluator may be correlated. Engineering assurance therefore cannot rest solely on the judgment that produced the result.

% <!-- index-def: long-horizon-task -->

\textbf{Bounded working state.} The information available to a foundation model during an invocation is finite. An agentic harness can extend its effective working state through retrieval, external memory, compaction, and repeated tool use, but information needed for a particular decision must be available when that decision is made. A \emph{long-horizon task} requires carrying earlier decisions or commitments into later steps. Large software systems exceed what can be presented in a single reasoning episode, so agents must select and reconstruct relevant system knowledge.

\textbf{Inexpensive repeated reasoning.} Agents can explore alternatives at a speed and cost that human engineers cannot match. An agent can propose an implementation, run checks, revise its work, and try again, often many times in the time a human engineer would need for one attempt. This is particularly valuable when the engineer can state what constitutes an acceptable result. Human engineers are fallible reasoners too: they misunderstand requirements, overlook interactions, make incorrect assumptions, and introduce defects. Software engineering has long responded with specifications, types, tests, reviews, and other means of checking work independently of its author. Artificial reasoners introduce a new source of implementation capacity, not a new need for engineering assurance.

% <!-- point: natural-language-delegation-needs-durable-intent | Natural-language requests can leave consequential choices implicit, and iterative changes can lose earlier obligations; FlowCheck provides a concrete scheduler example. | terms: modeling-principle, engineered-environment -->

These capabilities make natural language a productive interface for delegating implementation. They also make it easy to leave consequential decisions implicit. In the practice known as \emph{vibe coding}, developers accept and revise generated programs without closely examining their implementation. Later requests may inadvertently displace earlier requirements. In the FlowCheck study, three coding systems were asked to modify a speaker scheduler to allow multiple speakers in one time slot. All three lost an earlier rule limiting each speaker to one slot, although the new request did not require removing it~\citep{vir2026flowcheck}. The earlier obligation needed to survive the later instruction.

% <!-- point: the-reasoning-horizon-defined | The reasoning horizon is the practical limit of work the reasoner and its supporting machinery can carry through reliably, given the available representations and resources; task demands are distinct from this system-level limit. | terms: reasoning-horizon, context-window -->

% <!-- index-def: reasoning-horizon -->

% <!-- box-family: inset -->

% <!-- inset-domain: AGENT HARNESS -->

% mage-box-prefix: Inset —

\begin{bookbox}{inset}{Context Windows, Compaction, and the Reasoning Horizon}
\label{inset-context-windows-and-compaction}
A long-running engineering task can fail even when an agent has enough information for each step in isolation. An earlier decision or test result may be missing when the next step depends on it. The \textbf{reasoning horizon} is the practical limit of the work a reasoner and its supporting machinery can carry through reliably. In seqBench the term describes the length of a sequence a model can reliably complete~\citep{ramezanali2025seqbench}; here it includes what the harness and external resources make possible.

A \textbf{context window} limits how much information a foundation model can process in one invocation. Instructions, conversation history, files, and tool results compete for space measured in tokens. A harness can carry work beyond one window by storing records for retrieval or using \textbf{compaction}: replacing accumulated context with a shorter account. Pages of investigation and test output might become a handoff recording the plan, decisions, checks, and remaining work. If the handoff says only ``the tests passed,'' the next agent cannot tell which tests ran or which version they exercised. Keeping the original results helps only if the agent retrieves them when needed~\citep{rajasekaran2025contextengineering}.

The generation of commodity intelligence that motivated this book is built around language models with bounded context windows. We deliberately abstract from that machinery: the \textbf{reasoning horizon} concerns what a reasoner and its supporting machinery can carry through reliably, given the available representations and resources. The concept is thus broader than a context-window limit, and is intended to remain useful as new kinds of intelligence bring different capabilities and constraints.
\end{bookbox}

\section{Engineering the Environment}\label{heading:engineering-the-agent-environment}

% <!-- section-terms: agentic-harness, engineered-environment, injection-point -->

% <!-- point: three-harness-capabilities | The harness supplies action through tools, context management, and interposition; it may also support task management and parallel execution. | terms: agentic-harness, injection-point -->

% <!-- index-def: agentic-harness -->

% <!-- index-def: engineered-environment -->

The harness supplies facilities for autonomous work; the engineered environment determines how those facilities are used for a particular system. A harness manages context, invokes tools, and may provide memory, parallel execution, and hooks through which actions can be observed or intercepted. The environment includes the existing codebase and the knowledge embodied in it, together with the requirements, architecture, documentation, models, tests, permissions, and operational infrastructure through which the engineer understands and governs the system. The distinction is one of roles: a hook facility belongs to the harness; installing a project check through it is part of engineering the environment. For example, the hook may run a check that refuses edits to generated files. Build checks and review procedures also belong to the environment even when they operate independently of the harness.

% <!-- print-columns-start: reasoning-horizon-fig -->

Research on external representations has long shown that how information is organized affects the effort needed to reason with it~\citep{larkin1987diagram,zhang1994representations}. A useful engineering model preserves relationships that would otherwise have to be reconstructed from implementation. \autoref{reasoning-horizon-fig} illustrates this distinction: the agent can recover relevant structure from raw code or retrieve it from a purposeful representation. Better representations can reduce repeated reconstruction effort without necessarily increasing the range of tasks the system can handle.

\begin{figure}[htbp]
\centering
\includegraphics[width=\linewidth]{book/latex/assets/reasoning-horizon-7eba103f30.pdf}
\caption{\emph{Representation and the reasoning horizon.} When relevant system knowledge exceeds a reasoner's working state, it must be recovered for later decisions. Purposeful representations can reduce the work of recovering that knowledge. This may lower the cost of a task or make a larger task feasible.}
\label{reasoning-horizon-fig}
\end{figure}

Explicit engineering models are valuable not because they are inherently more accurate than a foundation model's internal representations, but because engineers can inspect, revise, and use them to enforce obligations. The environment can influence an agent's decisions by supplying knowledge and instructions; it can also constrain actions and evaluate results. Supplying an architectural model may help an agent preserve a dependency boundary, while a validator can detect a forbidden dependency and a build gate can refuse the change even when the agent fails to follow the model. These are different ways of exercising engineering control. MAGE develops them through Modeling and Alignment in the chapters that follow.

% <!-- print-columns-end: reasoning-horizon-fig -->

\section{A Simple Probabilistic Model of Success with Agents}\label{heading:a-simple-probabilistic-model-of-success-with-agents}

% <!-- section-terms: foundation-model, engineering-substrate, enforcement-problem, alignment-principle -->

Commodity intelligence is probabilistic. This book focuses on the implications for engineering. A foundation model's answer does not ordinarily come with a guarantee that its engineering judgments are correct. Generation can vary between runs, a request may admit several interpretations, and the model may choose different reasoning steps or tools. A task can depend on many such judgments before it produces a result. The reliability of the complete process also depends on how the harness and engineered environment organize the work.

% <!-- print-keep-with-next: 7 -->

A simple model illustrates why reliability on individual features does not guarantee reliability across a product. Suppose an agent implements \(n\) features, with \(p_i\) denoting the probability that feature \(i\) is implemented correctly. Under the simplifying assumption that these outcomes are independent, the probability that all features are correct is

\[
P(\text{all features correct})=\prod_{i=1}^{n}p_i
\]

When probabilities between zero and one are multiplied, their product decreases as more factors are added. If an agent implements each feature correctly with probability 99\%, the probability that all three features are correct is \(0.99^3 \approx 97\%\). For ten features it is approximately 90\%, and for one hundred it falls to approximately 37\%. At 95\% reliability per feature, the probability of implementing all ten correctly is approximately 60\%. An agent that usually gets individual features right can therefore have a much lower probability of getting every feature in a product right.

The independence assumption is deliberately simple; real implementation failures may be correlated, and a product's features are not necessarily equally difficult. The calculation is not a prediction of agent performance. It illustrates why the reliability of the whole system cannot be inferred from success on individual tasks. Human review, automated checks, iterative repair, and more capable models can improve the final result, but their costs and limitations belong in the evaluation of the complete engineering process.

% <!-- index-def: probabilistic-surface -->

% <!-- gloss-only: Probabilistic surface | The decisions in a system that depend on agentic inference. A larger probabilistic surface creates more opportunities for inference error to affect the outcome, even when each individual decision is highly reliable. MAGE seeks to reduce this surface by moving suitable decisions into explicit models, deterministic mechanisms, and independently enforced obligations, while leaving genuinely judgment-dependent decisions probabilistic. -->

% <!-- point: the-probabilistic-surface-is-the-judgments-still-bought-per-instance | Decisions that depend on agentic inference form the probabilistic surface; mechanisms can enforce particular requirements even when the agent makes a mistake. | terms: engineering-substrate, foundation-model -->

The decisions that depend on agentic inference form the system's \textbf{probabilistic surface}. Other parts may be deterministic, such as a tool that applies a fixed transformation or a test suite that evaluates explicit conditions. Better models and iterative loops can improve the reliability of inferred judgments. Engineers can also establish mechanisms that enforce selected requirements even when an agent makes a mistake. The engineering objective is not to eliminate probabilistic reasoning, but to decide where its uncertainty is acceptable and where stronger controls are needed.
